\honest{Belief in Belief}{Eliezer Yudkowsky}{2007}

Carl Sagan told a parable. Someone says, ``There is a dragon in my garage,'' and then that it is invisible. Such a claim could still be tested: we might hear breathing or measure carbon dioxide. But for each test we propose, the claimant has an excuse ready. The dragon is inaudible, does not breathe, and is ``permeable to flour.''

Sagan's moral is that poor hypotheses dodge falsification. Mine is different. The claimant knows in advance exactly which results need excusing, so somewhere in their mind they must have an accurate model: they expect no dragon. \nb{Sagan concluded that ``something funny is going on inside my head,'' but did not ask what the claimant expects. The inference from the advance excuses is the post's own.}

Some philosophers are ``much confused'' by this and ask whether the claimant really believes in the dragon, ``As if the human brain only had enough disk space to represent one belief at a time!'' The claimant does not anticipate seeing a dragon, or there would be no advance excuses; the sentence ``There is a dragon in my garage'' may still sit among the claimant's verbal beliefs. The two ought to clash, but you can write ``The sky is green!'' next to a picture of a blue sky and the paper does not burst into flames. \nb{Their question is the standard puzzle of self-deception: how can a person hold contradictory beliefs at once? Answers proposed before 2007 include Robert Audi's, that the self-deceived sincerely avow what they do not believe, which is close to the answer this post gives. The post names no philosopher.}

Asking what a belief predicts will not cure this, I admit. It is easy to connect the dragon to experience: if there is a dragon, you expect to see one when you open the door, and if you see none, there is none. The invisibility is a symptom of something worse.

Daniel Dennett calls the real problem ``belief in belief'': it is often easier to believe that you ought to believe something than to believe it. It is not even clear what you would believe if you believed that the Ultimate Cosmic Sky is perfectly blue and perfectly green, but it is easy to believe that believing it is good and virtuous. I think even Dennett oversimplifies. If you only believe in belief, you cannot admit it to yourself, since what is virtuous is to believe. ``Nobody will admit to themselves'' otherwise, unless unusually honest: people do not believe in belief in belief, they just believe in belief. (Mathematical logic trains you to keep P, a proof of P, and a proof that P is provable apart; the same sharp lines separate P, wanting P, believing P, wanting to believe P, and believing that you believe P.) I give no evidence for this beyond the parable and the reader's memory of doubting Santa Claus.

Belief in belief can be explicit, or it can be a fear of ridicule or a flinch that protects your self-image as the discoverer of the dragon. So we need ``a wider concept of belief, not limited to verbal sentences,'' including wordless expectations. \nb{Dennett had drawn a similar line himself in 1978, between beliefs, which animals also have, and verbal ``opinions,'' and had placed self-deception where the two come apart. This comes from Keith Frankish's account of Dennett; Dennett's chapter could not be checked here.} You can expect no dragon when you open your garage without ever thinking the sentence ``There is no dragon in my garage,'' and you can flinch from admitting you do not believe without ever thinking ``I want to believe.'' So it is not realistic to say the claimant believes it is beneficial to believe in the dragon; it is realistic to say the claimant ``anticipates as if'' there is no dragon and ``makes excuses as if they believed in the belief.''

A person who believes in the dragon and in their belief will risk a test. But excuses made in advance, ``it would seem,'' show that belief and belief in belief have come apart.
